\documentclass[11pt,a4paper]{extarticle}

\usepackage{kotex}
\usepackage[margin=17mm]{geometry}
\usepackage{amsmath,amssymb,mathtools}
\usepackage{array}
\usepackage{enumitem}
\usepackage{mdframed}
\usepackage{needspace}
\usepackage[hidelinks]{hyperref}

\setlength{\parindent}{0pt}
\setlength{\parskip}{5pt}
\setlist[itemize]{leftmargin=18pt,itemsep=2pt,topsep=2pt}
\setlist[enumerate]{leftmargin=22pt,itemsep=3pt,topsep=3pt}
\setcounter{tocdepth}{2}

\newcommand{\CM}{\mathrm{CM}}
\newcommand{\Poseidon}{\mathrm{Poseidon2}}
\newcommand{\Hash}{\mathrm{H}}
\newcommand{\HToF}{\mathrm{HashToField}}
\newcommand{\FrHash}{\mathrm{fr.Hash}}
\newcommand{\Verify}{\mathrm{Verify}}
\newcommand{\Prove}{\mathrm{Prove}}
\newcommand{\Append}{\mathrm{Append}}
\newcommand{\Member}{\mathrm{Member}}
\newcommand{\Range}{\mathrm{Range}}
\newcommand{\sample}{\xleftarrow{\$}}
\newcommand{\DocumentHash}{\mathrm{DocumentHash}}
\newcommand{\Document}{\mathrm{Document}}
\newcommand{\addrSender}{addr_{Sender}}
\newcommand{\addrReceiver}{addr_{Receiver}}
\newcommand{\DeltaEpoch}{\Delta_{epoch}}
\newcommand{\Assert}{\mathbf{assert}\;}
\newcommand{\Return}{\mathrm{return}\;}
\newcommand{\List}[1]{\mathrm{List}_{#1}}
\newcommand{\Map}[1]{\mathrm{Map}_{#1}}
\newcommand{\MT}{\mathrm{MT}}
\newcommand{\rvMT}{\mathrm{rvMT}}
\newcommand{\CurrentEpoch}{\mathrm{CurrentEpoch}}

\newenvironment{schemeblock}
{%
  \begingroup
  \setlength{\jot}{8pt}%
  \setlength{\parskip}{0pt}%
  \setlength{\abovedisplayskip}{9pt}%
  \setlength{\belowdisplayskip}{0pt}%
  \setlength{\abovedisplayshortskip}{9pt}%
  \setlength{\belowdisplayshortskip}{0pt}%
  \begin{mdframed}[
    linewidth=0.4pt,
    linecolor=black,
    innerleftmargin=8pt,
    innerrightmargin=8pt,
    innertopmargin=8pt,
    innerbottommargin=8pt,
    skipabove=0pt,
    skipbelow=5pt,
    splittopskip=8pt,
    splitbottomskip=8pt
  ]
  \centering
}
{%
  \end{mdframed}
  \endgroup
}

\title{zkDPP Protocol Specification}
\author{Canonical Protocol Specification}
\date{2026-07-16}
\renewcommand{\contentsname}{목차}

\begin{document}
\maketitle
\tableofcontents
\newpage

\section{프로토콜 기준과 공통 객체}

\subsection{문서 역할}

이 문서는 zkDPP의 canonical protocol semantics를 정의한다. Protocol object,
Client와 Smart Contract의 역할, public/private statement, 8개 Relation, contract 검사와
상태 갱신 경계를 이 문서의 source of truth로 사용한다. 특정 library API, package type,
serialization code와 test harness는 정의하지 않는다. 해당 구현 선택은
\texttt{Dev Spec.tex}에서 정의한다.

\subsection{프로토콜 개요}

zkDPP는 product information을 숨긴 finalized note를 공급망 Event 사이에서 소비하고
새로 생성하면서, quantity와 consistency-critical StateVector의 전이가 각 Event
Relation을 만족했음을 증명한다.

\begin{center}
\begin{tabular}{lll}
\hline
구분 & Event & Protocol object 전이 \\
\hline
공급망 & Entry & $\varnothing\rightarrow cm_{new}$ \\
공급망 & Transfer & $cm_{old}\rightarrow (cm_{change},rv_{new})$ \\
Resolution & Proceed & $rv_{old}\rightarrow cm_{new}$ \\
Resolution & Recall & $rv_{old}\rightarrow cm_{new}$ \\
공급망 & Merge & $(cm_{old,1},cm_{old,2})\rightarrow cm_{new}$ \\
공급망 & Split & $cm_{old}\rightarrow (cm_{new,1},cm_{new,2})$ \\
공급망 & Process & $(cm_{old,i})_{i=1}^{m}\rightarrow(cm_{new,j})_{j=1}^{n}$ \\
공급망 & Exit & $cm_{old}\rightarrow\varnothing$ \\
\hline
\end{tabular}
\end{center}

Finalized $cm$은 $\MT$에, pending Recall Voucher $rv$는 $\rvMT$에 append한다.
$nf$는 finalized note 소비를, $rvnf$는 Recall Voucher의 단일 resolution을 강제한다.

\subsection{암호학적 추상화}

각 Relation $\mathcal{R}_i$에는 proving key와 verifying key가 존재한다.
$\mathrm{Setup}_{Client}$는 Relation 집합을 입력받아 공통 system parameter와 각
Relation의 key pair를 생성한다. $pp$에는 system parameter와 proof key가 함께 들어가지만,
$\mathrm{KeyGenUser}_{Client}$는 이 중 Field와 hash 같은 system parameter만 사용하고
PLONK proving/verifying key는 사용하지 않는다. Proof system, curve, library version,
POC SRS 생성 방식과 verifier 배포 방식은 Development Specification에서 고정한다.

$\Hash$는 key derivation, nullifier와 Recall Voucher derivation에 사용하는
추상 hash interface다. POC에서는 Poseidon2를 사용하지만, Scheme에는 package API나
parameter object를 노출하지 않는다. $nf=\Hash(sk_{own},cm)$은 기존
PRF 기반 nullifier가 담당하는 deterministic secret-key derivation 역할을 hash로
구현한다는 수준으로 해석한다. Hash가 PRF와 동일하다는 일반적 주장을 의미하지는 않는다.

$\CM$은 finalized note를 binding하는 추상 commitment constructor다. POC에서는
Poseidon2로 구현하지만, Scheme은 입력값과 binding relation만 정의한다.

\subsection{확정 Protocol 경계}

\begin{itemize}
  \item Circuit 집합은 Entry, Transfer, Proceed, Recall, Merge, Split,
        Process, Exit의 8개다.
  \item Merge와 Split은 각각 fixed 2-to-1, 1-to-2 circuit으로 구현한다.
  \item Process는 fixed maximum input/output arity $M_{\max}$와 $N_{\max}$를 갖는다.
        기능 확인용 기본값은 $M_{\max}=N_{\max}=3$이다. 성능 측정은
        $M_{\max}=N_{\max}=a$, $a\in\{1,2,3\}$인 별도 circuit을 사용한다.
  \item 각 Process profile은 $1\leq m,n\leq a$를 지원한다. $m,n$은 active input/output
        개수를 나타내는 public Field Element이며 활성 slot은 index 1부터 연속된다.
        별도 public 또는 private enable vector는 두지 않는다. Inactive input의 $rt$, $cm$,
        $nf$와 inactive output의 $cm$, $A$ row는 0이다.
  \item Process의 semantic delta는 State prefix에 따라
        $P=(-p_{totalLoss},+p_{carbonAdded},+p_{recycledDelta})$의 앞 $\ell$개
        coordinate를 사용한다.
        Public input은 각 State coordinate에 대응하는 unsigned 값을 포함한다.
        $\ell=3$에서는
        $P_{pub}=(p_{totalLoss},p_{carbonAdded},p_{recycledDelta})$이고
        recycled delta 적용식과 recycled aggregate 보존식을 함께 검증한다. 따라서
        유효한 proof에서는 $p_{recycledDelta}=0$이 도출되며 별도 0 assertion은 두지 않는다.
        Circuit은 $\mathbf{s}_{in}$을 input State aggregate,
        $\mathbf{s}_{out}$을 output State aggregate로 계산하고
        $\mathbf{s}_{out}=\mathbf{s}_{in}+P$에 대응하는 coordinate별 relation을 검증한다.
  \item Process allocation denominator는 circuit constant $D=10^9$으로 고정하고,
        StateAllocationMatrix $A$는 public input으로 둔다.
  \item $q$는 unsigned 64-bit integer로 제한한다. Transfer의
        $cm_{change}$는 $q=0$인 경우에도 항상 생성한다.
  \item StateVector의 각 coordinate는 unsigned 64-bit integer이며,
        nano scale $10^9$ fixed-point 값으로 encode한다.
  \item Entry는 $0<q<2^{64}$와
        $0\leq s_k<2^{64}$를 검사한다. Entry의 신규 Product quantity는 0일 수 없지만,
        State coordinate는 0을 허용한다. 별도의 initial product semantic invariant는
        강제하지 않는다.
  \item Transfer의 output ordering은 $rv$를 첫 번째, $cm_{change}$를 두 번째로 둔다.
        Transfer와 Split은 deterministic first-output remainder를 사용한다.
        Process도 output 1이 모든 State coordinate의 residual을 받는다.
  \item Field 내부의 aggregate, cross multiplication, Process allocation에서
        modular wraparound가 발생하지 않도록 필요한 intermediate bound는
        circuit별 range 분석 후 확정한다.
  \item Pedersen commitment는 사용하지 않는다. Note opening과 event relation은
        circuit 내부에서 Poseidon2 commitment에 직접 binding한다.
  \item POC의 $\Hash$와 $\CM$은 gnark v0.15.0과 gnark-crypto v0.20.1이 제공하는
        BLS12-381 default Poseidon2 hasher를 사용한다. 구체 parameter와 reference
        vector는 Development Specification에서 고정한다.
  \item 이후 POC 구현과 benchmark는 Poseidon2 profile만 사용한다. SHA-to-Field
        profile은 EVM SHA-256 precompile의 gas와 circuit 비용을 확인하기 위한 일회성
        비교로 종료했으며, 이미 얻은 결과와 재현 코드는 Development Specification에
        보존한다. 새로운 SHA circuit, setup, proof와 gas 실험은 수행하지 않는다.
  \item Azeroth의 BLS12-381 Solidity 파일은 protocol/verifier separation과
        Merkle hash abstraction만 참고한다. Groth16 verifier, 과거 EIP-2537 ABI와
        BN254용 Poseidon2/MiMC source는 재사용하지 않는다.
  \item DocumentInfo와 DocumentHash의 일치 여부는 circuit 또는 contract에서
        검증하지 않는다. Receiver가 off-chain에서 재계산한다.
  \item Input/output commitment edge를 공개하는 provenance event는 현재 구현
        범위에서 제외한다. 다만 Client가 Merkle witness를 구성할 수 있도록
        $MT$와 $rvMT$ append의 leaf와 zero-based index는 event로 공개한다.
\end{itemize}

\subsection{POC 구현 Profile}

Protocol의 일반 모델은 위 경계를 유지한다. 최초 vertical slice는 Entry와 고정
3-to-3 Process로 시작했으며, 현재 \path{poc-v2/}는 8개 Event를 모두 구현한다.
Canonical 기능 profile은 다음과 같다.

\begin{itemize}
  \item StateVector 길이 $\ell=3$, Merkle Tree depth $d=32$
  \item Entry, Transfer, Proceed, Recall, Merge, Split, Process와 Exit를 하나의
        canonical scenario에서 순서대로 실행
  \item Process 기능 기준은 정확히 3 input과 3 output이며 public $m,n$은 모두 3
  \item Go proof와 generated Solidity verifier를 함께 검증하고 Main Contract의
        $MT$, $rvMT$, nullifier와 voucher state update까지 확인
\end{itemize}

성능 profile은 State 길이 0부터 3, Process arity 1-to-1부터 3-to-3을 별도 circuit으로
compile한다. 이는 일반 Protocol relation을 변경하지 않는 benchmark 전용 compile-time
specialization이다.

\subsection{관련 문서}

Data encoding, hash parameter, Merkle tree, proof serialization, contract ABI와 test
세부사항은 같은 폴더의 \texttt{Dev Spec.tex}을 따른다.
결정 이유와 반려된 선택지는 \texttt{Decision Log.tex}에 기록한다.

\subsection{문서 구성}

1장은 zkDPP Scheme을 읽기 위한 공통 객체와 표기를 설명한다. 2장은 Client와
Smart Contract의 초기화를 정의한다. 3장은 Azeroth Scheme과 같은 형식으로 각
Event를 Client 함수, circuit relation, Smart Contract 함수 순서로 배치한다.

공급망 event는 Entry, Transfer, Merge, Split, Process, Exit이다. Proceed와 Recall은
Transfer가 생성한 Recall Voucher를 해소하는 resolution event다. 문서의 event 순서는
Entry, Transfer, Proceed, Recall, Merge, Split, Process, Exit이다.

\subsection{DocumentInfo와 DocumentHash}

\[
  \mathit{DocumentInfo}=(ProductName,LotId,Unit)
\]

\[
  b_x=\operatorname{UTF8}(\operatorname{NFC}(x))
\]

\[
  \operatorname{EncodeText}(x)
  =
  \operatorname{U32BE}(|b_x|)
  \mathbin{\|}b_x
\]

\[
  \begin{aligned}
  \mathit{DocumentInfoBytes}
  ={}&\operatorname{EncodeText}(ProductName)\\
     &\mathbin{\|}\operatorname{EncodeText}(LotId)\\
     &\mathbin{\|}\operatorname{EncodeText}(Unit)
  \end{aligned}
\]

\[
  \DocumentHash
  =
  \FrHash(
    \mathit{DocumentInfoBytes},
    \varepsilon,1
  )[0]
\]

$\mathit{DocumentInfo}$는 circuit input이 아니다.

POC의 \texttt{HashToField}는 BLS12-381 scalar field의 \texttt{fr.Hash}를
사용한다. 이 함수는 SHA-256 기반 XMD로 byte string을 확장한 뒤
$\mathbb{F}_r$ 원소로 변환한다. 두 번째 인자는 API에서 요구하는 DST지만,
POC는 별도의 hash 용도 구분값을 사용하지 않고 빈 byte string $\varepsilon$을
전달한다. 직렬화에는 layout version을 넣지 않는다. 각 문자열은 Unicode NFC로
정규화한 뒤 UTF-8로 encode한다. 해당 UTF-8 byte 길이를 4-byte unsigned
big-endian으로 먼저 기록하고 UTF-8 bytes를 이어 붙인다. 수령자는 off-chain에서 위의
$\DocumentHash$ 식을 재계산한다.
Circuit은 이 hash가 event 사이에서 유지되거나 정해진 relation에 따라
변경되는지만 검증한다.

\subsection{Owner, StateVector, finalized note}

별도 설명이 없는 한 이 문서의 Field는 BLS12-381 scalar field
$\mathbb{F}_r$를 뜻한다. Base field $\mathbb{F}_q$는 curve point의 좌표와
pairing 구현에 사용하며 protocol payload를 담는 Circuit Field가 아니다.

\[
  sk_{own},pk_{own},addr,\DocumentHash,q,s_k,o,cm
  \in\mathbb{F}_r
\]

\[
  pk_{own}=\Hash(sk_{own})
  \,;\,
  addr=\Hash(pk_{own})
\]

\[
  q\in\mathbb{Z}_{\geq0}
  \,;\,
  \mathbf{s}=(s_1,\ldots,s_{\ell})
\]

\[
  k\in\{1,\ldots,\ell\}
  \,;\,
  (s_1,s_2,s_3)=(total_{kg},carbon_{kgCO2e},recycled_{kg})
\]

기능 확인용 기본 POC는 $\ell=3$을 사용한다. 성능 확장성 측정에서는
$\ell\in\{0,1,2,3\}$인 별도 circuit profile을 compile하며, 각 profile의 StateVector는
위 순서의 prefix를 사용한다. 따라서 $\ell=0$은 empty StateVector, $\ell=1$은
$total_{kg}$, $\ell=2$는 $total_{kg}$와 $carbon_{kgCO2e}$, $\ell=3$은 세 coordinate
전체를 뜻한다. Profile 사이에서 note, proof와 verifier를 혼용하지 않는다.

\[
  \begin{aligned}
  \CM(\DocumentHash,q,\mathbf{s},addr,o)
  =\Poseidon(\DocumentHash,q,s_1,\ldots,s_{\ell},addr,o)
  \end{aligned}
\]

\[
  \mathit{note}=(cm,\DocumentHash,q,\mathbf{s},o)
\]

\texttt{quantity}는 integer다. StateVector coordinate는 nano scale $10^9$의
unsigned fixed-point integer다. $q$와 $s_k$는 Field Element이면서 동시에
unsigned 64-bit range를 만족해야 한다. $pk_{own}$은 elliptic-curve point가 아니라
$sk_{own}$에서 Poseidon2로 도출한 public ownership identifier다. Note는 $addr$을
저장하지 않는다. Holder는
$\mathrm{KeyGenUser}_{Client}$에서 받은 $addr$을 별도로 알고 있으며, spend circuit은
$addr=\Hash(\Hash(sk_{own}))$를 검증한다.

\subsection{Nullifier, Recall Voucher, on-chain state}

\[
  nf=\Hash(sk_{own},cm)
\]

\[
  rv=
  \Hash(
    \DocumentHash,q_T,\mathbf{s}_T,
    \addrSender,\addrReceiver,\DeltaEpoch,o_{rv}
  )
\]

\[
  rvnf=\Hash(rv,o_{rv})
\]

\[
  \mathit{voucher}=\{
    rv,\DocumentHash,q_T,\mathbf{s}_T,
    \addrSender,\addrReceiver,\DeltaEpoch,o_{rv}
  \}
\]

$nf$는 finalized note의 중복 소비를 방지한다. $rvnf$는 하나의 Recall Voucher가
Proceed와 Recall 중 정확히 한 번만 해소되도록 한다. $o_{rv}$는 sender와 receiver가
off-chain으로 공유한다. $\mathit{voucher}$는 Sender와 Receiver가 동일한 schema로
보관하는 canonical off-chain object다. 각 holder가 자신의 address를 이미 알고 있어도
역할별 schema를 나누지 않고 두 address를 모두 저장한다.

Recall Voucher의 resolution 상태는 raw $rv$를 key로 하는 $resolved[rv]$ mapping이 아니라
$rvnf=\Hash(rv,o_{rv})$와 $\List{rvnf}$로 관리한다. Proceed와 Recall transaction에는
$rv$와 $rvnf$가 함께 공개되므로 이 선택이 둘 사이의 on-chain linking을 숨긴다고
주장하지 않는다. 목적은 finalized note의 $nf$와 일관된 secret-derived nullifier 구조와
Sender/Receiver가 동일한 resolution identifier를 계산하도록 하는 것이다.

Finalized note와 Recall Voucher는 서로 다른 append-only Merkle tree에 저장한다.

\[
\operatorname{depth}(MT)=\operatorname{depth}(rvMT)=d,
\qquad
d=32\ \text{for the POC}
\]

Depth scaling은 수행하지 않는다. State scaling과 Process arity scaling은
$d=32$로 고정한다. 별도 Merkle microbenchmark에서 depth-32 commitment 한 개의
membership verification 비용만 측정한다.

\[
z_0=0,
\qquad
z_{h+1}=H_{\mathrm{tree}}(z_h,z_h),
\qquad
root_{\mathrm{empty}}=z_d
\]

Width-2 default Poseidon2 permutation을 $\mathcal{P}_{P2}$라 할 때 Tree compressor는
다음으로 고정한다.

\[
(u_0,u_1)=\mathcal{P}_{P2}(left,right),
\qquad
H_{\mathrm{tree}}(left,right)=u_1+right\pmod r
\]

즉 초기 두 lane은 $(left,right)$이고, permutation 이후 오른쪽 lane에 원래
$right$를 feed-forward한 값을 출력한다. Native Go와 gnark circuit은 pinned library의
동일한 \texttt{Compress(left,right)} 의미론을 사용한다.

$MT$에는 $cm$을, $rvMT$에는 $rv$를 leaf-hash 없이 그대로 append한다. Valid $cm$과
$rv$는 0이 아니지만, Azeroth와 동일하게 별도 nonzero Contract 검사는 하지 않는다.
두 Tree는 왼쪽부터 순차 append하며 각 Tree의 empty root와 모든 과거 root를 accepted
root로 유지한다. Membership witness는 leaf index와 leaf-to-root 순서의 sibling
$d$개로 구성한다. Append 후에는 leaf, zero-based index와 new root를 event로
공개한다. 구체 storage layout과 $H_{\mathrm{tree}}$ 구현은 Development
Specification을 따른다.

\[
  \List{rt},\List{cm},\List{nf},
  \List{rvrt},\List{rv},\List{rvnf},
  \Map{user},\Map{rv}
\]

POC의 8개 Event 함수는 permissionless다. Transaction을 제출한 Ethereum
\texttt{msg.sender}를 note owner로 요구하지 않으며, Note와 Recall
Voucher의 소비 권한은 Circuit relation으로 검증한다. User registration은 호출 EOA를
등록한 $(addr,pk_{own})$에 연결한다. Logical List는
membership set이고, Solidity의 구체 mapping과 Tree storage layout은 Development
Specification을 따른다.

각 Smart Contract Event 함수는 원칙적으로 다음 세 단계를 따른다.

\begin{enumerate}
  \item public argument, accepted root, 기존 List 중복, nullifier와 deadline 같은
        on-chain precondition을 검사한다.
  \item ordered public statement를 구성하고 해당 Relation proof를 검증한다.
  \item proof가 유효할 때만 List와 Map을 갱신하고 $MT$ 또는 $rvMT$에 append한다.
\end{enumerate}

Proceed와 Recall은 accepted $rvMT$ root와 Circuit membership proof에 더해
$rv\in\List{rv}$도 사전 검사한다. 이 검사는 membership soundness를 대신하지 않으며,
명백히 등록되지 않은 $rv$를 proof 검증 전에 거부하는 저비용 on-chain filter로 유지한다.

다중 output의 같은 transaction 내부 중복에는 pairwise 비교를 추가하지 않는다. Proof
검증 후 output을 순서대로 기존 List에 없는지 검사하고 등록한 다음 Tree에 append한다.
뒤 output이 앞 output과 같으면 순차 check-and-insert가 실패하며, transaction atomicity에
따라 해당 transaction의 앞선 List, Tree와 event 변경도 모두 revert된다.

\[
  EPOCH\_SIZE=600\ \mathrm{seconds}
\]

\[
  \CurrentEpoch()
  =
  \left\lfloor\frac{block.timestamp}{EPOCH\_SIZE}\right\rfloor
\]

\[
\Map{rv}[rv]
  =
  \CurrentEpoch()+\DeltaEpoch
\]

Sender는 Transfer를 생성할 때 $\Delta_{epoch}$을 선택한다. Protocol policy는
$\Delta_{epoch}\geq0$만 요구하며 별도의 business maximum을 두지 않는다.
$\Delta_{epoch}$은 public Field Element로 voucher에 binding되지만 Circuit에서 별도
bit range check를 하지 않는다. Generated verifier가 모든 public input과 마찬가지로
$0\leq\Delta_{epoch}<r$인 canonical scalar-field 표현인지 확인하므로 Main Contract는
같은 검사를 반복하지 않는다.

Transfer가 실행될 때 Smart Contract가 현재 epoch를 계산하고, 여기에
$\Delta_{epoch}$을 더한 값을 $\Map{rv}[rv]$에 저장한다. 외부에서 완성된 deadline
값을 입력받아 저장하지 않는다. Deadline은 proof 검증 후 일반 \texttt{uint256} 덧셈으로
계산한다. $\Delta_{epoch}<r<2^{255}$이고
$\CurrentEpoch()<2^{256}/600<2^{247}$이므로 protocol input 범위에서 합은
\texttt{uint256}을 overflow하지 않는다. Recall 또는 Proceed 이후에도 이 entry를
삭제하지 않는다.

Recall은 $\CurrentEpoch()<\Map{rv}[rv]$일 때만 허용한다. 따라서
$\Delta_{epoch}=0$인 Transfer는 생성 직후부터 Recall할 수 없다. Proceed는 voucher가
아직 resolve되지 않았다면 deadline 이후에도 허용한다.

Client는 leaf의 zero-based index를 Tree별 view 함수에 전달해 현재 root와 Merkle path를
조회하고 이를 witness로 준비한다.

\[
\mathrm{GetPath}_{MT}(index)\to(rt_{current},sibling_0,\ldots,sibling_{d-1})
\]

\[
\mathrm{GetPath}_{rvMT}(index)\to(rvrt_{current},sibling_0,\ldots,sibling_{d-1})
\]

각 spend Relation은 $\mathrm{Member}(rt,leaf,path)=1$을 Circuit 내부에서 검증한다.
Smart Contract는 path를 다시 계산하지 않고 public input $rt$가 accepted-root 목록에
있는지와 proof가 유효한지를 검사한다. View 함수가 반환하는 path는 조회 시점의 현재
root에 대한 path다. 이미 저장한 과거 root용 path는 Client가 별도로 보관할 수 있다.
새 leaf append와 새 root 계산은 Smart Contract가 수행한다.

2장 이후에서 $\mathbf{x}$는 public statement, $\mathbf{w}$는 private witness를 뜻한다.
$\mathbf{x}$와 $\mathbf{w}$는 중괄호로 표기하지만 unordered set이 아니라 ordered tuple이다.
각 $\mathbf{x}$ tuple에 적힌 왼쪽에서 오른쪽 순서가 Circuit과 Solidity verifier가
사용하는 canonical public-input 순서다.
Process는 on-chain recipe registry를 가정하지 않는다. 공개된 $A$는 해당 proof가
사용한 allocation을 나타내지만, $A$ 자체가 승인된 물리 공정 규칙임을 보장하지 않는다.
공개된 $P,A$는 선언된 State transition의 transparency와 accounting consistency를
제공할 뿐, 실제 물리 공정과의 일치나 산업적으로 정당한 allocation을 증명하지 않는다.
초기 State의 현실 정확성은 Entry issuer에, $P,A$의 현실 정확성은 Process operator 또는
외부 검증 주체에 의존한다. Recipe/Policy 기반 Material 및 Quantity transition 검증은
이번 POC 범위에서 제외한다.
Process는 input과 output의 DocumentHash 사이에 equality 또는 transition policy를
두지 않는다. 서로 다른 material과 product를 나타내므로 각 DocumentHash가 해당
input/output commitment에 binding되는지만 검증한다.

\section{초기화}

\subsection{\texorpdfstring{$\mathrm{Setup}_{Client} / \mathrm{Setup}_{SC}$}{Setup\_Client / Setup\_SC}}

{
\setlength{\jot}{8pt}
\setlength{\abovedisplayskip}{9pt}
\setlength{\belowdisplayskip}{0pt}
\setlength{\fboxsep}{8pt}

\noindent
\fbox{%
\begin{minipage}{0.96\textwidth}
\begin{tabular}{@{}p{0.475\linewidth}|p{0.475\linewidth}@{}}
\begin{minipage}[t]{\linewidth}
\centering
\(
  \mathbf{Setup}_{Client}
  (\{\mathcal{R}_i\}_{i\in\mathcal{I}})
\)

\begin{gather*}
\mathcal{I}=
  \{\mathrm{Entry},\mathrm{Transfer},\mathrm{Proceed},\mathrm{Recall}, \\
  \mathrm{Merge},\mathrm{Split},\mathrm{Process},\mathrm{Exit}\} \\
(ek_i,vk_i)
  \gets
  \mathrm{Setup}_{PLONK}(\mathcal{R}_i),
  \quad i\in\mathcal{I} \\
pp_{sys}\gets(\mathbb{F},\Hash,\CM,\Hash_{tree}) \\
pp\gets(pp_{sys},\{ek_i,vk_i\}_{i\in\mathcal{I}}) \\
\Return pp
\end{gather*}
\end{minipage}
&
\begin{minipage}[t]{\linewidth}
\centering
\(
  \mathbf{Setup}_{SC}
  (\{vk_i\}_{i\in\mathcal{I}})
\)

\begin{gather*}
vk_{SC,i}\gets vk_i,
  \quad i\in\mathcal{I} \\
\mathrm{Init}_{\MT}
  \,;\,
  \mathrm{Init}_{\rvMT} \\
\mathrm{Init}\bigl(
  \List{rt},\List{cm},\List{nf}, \\
  \List{rvrt},\List{rv},\List{rvnf}, \\
  \Map{user}, \Map{rv}
\bigr)
\end{gather*}
\end{minipage}
\\
\end{tabular}
\end{minipage}%
  }
}

\subsection{\texorpdfstring{$\mathrm{KeyGenUser}_{Client} / \mathrm{RegisterUser}_{SC}$}{KeyGenUser\_Client / RegisterUser\_SC}}

{
\setlength{\jot}{8pt}
\setlength{\abovedisplayskip}{9pt}
\setlength{\belowdisplayskip}{0pt}
\setlength{\fboxsep}{8pt}

\noindent
\fbox{%
\begin{minipage}{0.96\textwidth}
\begin{tabular}{@{}p{0.475\linewidth}|p{0.475\linewidth}@{}}
\begin{minipage}[t]{\linewidth}
\centering
\(
  \mathbf{KeyGenUser}_{Client}(pp_{sys})
\)

\begin{gather*}
sk_{own}\sample\mathbb{F}
  \,;\,
  pk_{own}\gets\Hash(sk_{own}) \\
addr\gets\Hash(pk_{own}) \\
usk=sk_{own}
  \,;\,
  upk=(addr,pk_{own}) \\
\mathit{Tx}_{KGU}=(addr,pk_{own}) \\
\Return (usk,upk),\mathit{Tx}_{KGU}
\end{gather*}
\end{minipage}
&
\begin{minipage}[t]{\linewidth}
\centering
\(
  \mathbf{RegisterUser}_{SC}(addr,pk_{own})
\)

\begin{gather*}
eoa\gets\mathrm{msg.sender} \\
\Map{user}[eoa]\gets(addr,pk_{own})
\end{gather*}
\end{minipage}
\\
\end{tabular}
\end{minipage}%
}
}

$\Map{user}$는 POC Client가 EOA로 $(addr,pk_{own})$을 조회하기 위한 편의용 metadata다.
Registration은 Event authorization에 사용하지 않으며 address uniqueness, EOA 재등록과
$addr=\Hash(pk_{own})$을 Contract에서 강제하지 않는다. $pk_{own}=\Hash(sk_{own})$과
$addr=\Hash(pk_{own})$에 의한 note ownership은 각 Event circuit에서 검증한다.

\section{공급망 이벤트}

각 Event 상단의 전이식은 소비되는 protocol object와 생성되는 protocol object를
요약한다. $cm_{old}$는 소비 대상, $cm_{new}$는 신규 생성 대상을 뜻한다. Transfer의
출력 ordered pair는 Recall Voucher multiset과 finalized $cm$ multiset 순서다. 세부
Client와 Relation에서는 역할에 따라 기존 $cm_{in}$, $cm_{out}$ 표기를 사용한다.

\subsection{Entry}

\[
  \varnothing
  \xrightarrow{\mathrm{Entry}}
  \{cm_{new}\}
\]

\noindent\textbf{표기.}
\begin{itemize}[nosep]
  \item $new$: 새로 발행하는 finalized note
\end{itemize}

Entry는 외부의 Product information을 새로운 finalized note로 발행하여 공급망의
$MT$에 등록하는 event다. POC는 issuer authorization, origin data의 사실성 또는
DocumentHash별 최초 발행의 유일성을 강제하지 않는다. 동일한 DocumentHash를 사용하더라도
fresh opening $o$와 note payload에 의해 서로 다른 $cm$을 생성할 수 있다.

Entry Relation은 $0<q<2^{64}$와 각 State coordinate의
$0\leq s_k<2^{64}$를 검사한다. 따라서 Entry quantity는 0일 수 없지만 State
coordinate는 0을 허용한다.

\subsubsection{\texorpdfstring{$\mathrm{Entry}_{Client}$}{Entry\_Client}}

\begin{schemeblock}
\(
  \mathbf{Entry}_{Client}
  (\DocumentHash,q,\mathbf{s},addr)
\)

\begin{gather*}
  o\sample\mathbb{F}; \, cm_{new}
  \gets
  \CM(\DocumentHash,q,\mathbf{s},addr,o)
  \\
  \mathbf{x}=\{cm_{new}\}
  \\
  \mathbf{w}=\{\DocumentHash,q,\mathbf{s},addr,o\}
  \\
  \pi_{Entry}
  \gets
  \Prove(ek_{Entry},\mathbf{x};\mathbf{w})
  \\
  \mathit{note}=(cm_{new},\DocumentHash,q,\mathbf{s},o)
  \\
  \mathit{Tx}_{Entry}=(\pi_{Entry},cm_{new})
  \\
  \Return (\mathit{note},\mathit{Tx}_{Entry})
\end{gather*}
\end{schemeblock}

\subsubsection{\texorpdfstring{$\mathcal{R}_{Entry}$}{R\_Entry}}

\begin{schemeblock}
\(
  \mathcal{R}_{Entry}(\mathbf{x};\mathbf{w})
\)

\begin{gather*}
  \Assert
  cm_{new}
  =
  \CM(\DocumentHash,q,\mathbf{s},addr,o)
  \\
  \Assert 0<q<2^{64}
  \,;\,
  \Assert 0\leq s_k<2^{64}
  \quad k\in\{1,\ldots,\ell\}
\end{gather*}
\end{schemeblock}

\subsubsection{\texorpdfstring{$\mathrm{Entry}_{SC}$}{Entry\_SC}}

\begin{schemeblock}
\(
  \mathbf{Entry}_{SC}(\pi_{Entry},cm_{new})
\)

\begin{gather*}
  \mathbf{x}=\{cm_{new}\}
  \\
  \Assert cm_{new}\notin\List{cm}
  \\
  \Assert
  \Verify(vk_{SC,Entry},\pi_{Entry},\mathbf{x})=1
  \\
  rt_{new}\gets\Append(\MT,cm_{new})
  \\
  \List{rt}.append(rt_{new})
  \,;\,
  \List{cm}.append(cm_{new})
\end{gather*}
\end{schemeblock}

\subsection{Transfer}

\[
  \{cm_{old}\}
  \xrightarrow{\mathrm{Transfer}}
  (\{rv_{new}\},\{cm_{new,change}\})
\]

\noindent\textbf{표기.}
\begin{itemize}[nosep]
  \item $in$: 소비하는 input note
  \item $T$: Receiver에게 transfer하는 pending portion
  \item $C$: Sender에게 남는 change portion
  \item $S$, $R$: 각각 Sender와 Receiver
  \item $k$: State coordinate index
  \item $r_{C,k}$: change State를 내림 계산할 때 생기는 나머지. Circuit witness로만
        사용하며 note, voucher와 transaction에는 저장하지 않는다.
\end{itemize}

\subsubsection{\textsc{TransferClient}}

\begin{schemeblock}
\(
  \textsc{TransferClient}(
    \mathit{note}_{in},usk_S,upk_S,upk_R,
    q_T,\DeltaEpoch,rt,path_{in})
\)

\begin{gather*}
(cm_{in},\DocumentHash,q_{in},\mathbf{s}_{in},o_{in})
  \Leftarrow
  \mathit{note}_{in}
  \\
sk_S\Leftarrow usk_S
  \,;\,
  (\addrSender,pk_S)\Leftarrow upk_S
  \,;\,
  (\addrReceiver,pk_R)\Leftarrow upk_R
  \\
q_C\gets q_{in}-q_T
  \\
s_{C,k}\gets
\left\lfloor\frac{q_Cs_{in,k}}{q_{in}}\right\rfloor
  \,;\,
r_{C,k}\gets q_Cs_{in,k}-q_{in}s_{C,k}
  \quad k\in\{1,\ldots,\ell\}
  \\
\mathbf{s}_T\gets\mathbf{s}_{in}-\mathbf{s}_C
  \\
o_C,o_{rv}\sample\mathbb{F}
  \\
nf\gets\Hash(sk_S,cm_{in})
  \\
rv\gets
  \Hash(
    \DocumentHash,q_T,\mathbf{s}_T,
    \addrSender,\addrReceiver,\DeltaEpoch,o_{rv}
  )
  \\
cm_C\gets
  \CM(\DocumentHash,q_C,\mathbf{s}_C,\addrSender,o_C)
  \\
\mathbf{x}_T=\{rt,cm_{in},nf,rv,cm_C,\DeltaEpoch\}
  \\
\begin{aligned}
  \mathbf{w}_T=\{
    &sk_S,pk_S,pk_R,\\
    &\DocumentHash,q_{in},\mathbf{s}_{in},
    q_T,\mathbf{s}_T,\addrSender,\addrReceiver,\DeltaEpoch,\\
    &q_C,\mathbf{s}_C,\mathbf{r}_C,path_{in},
    o_{in},o_C,o_{rv}\}
  \end{aligned}
  \\
\pi_T\gets
  \Prove(ek_{Transfer},\mathbf{x}_T;\mathbf{w}_T)
  \\
\mathit{note}_C=(cm_C,\DocumentHash,q_C,\mathbf{s}_C,o_C)
  \\
\mathit{voucher}=\{
    rv,\DocumentHash,q_T,\mathbf{s}_T,
    \addrSender,\addrReceiver,\DeltaEpoch,o_{rv}\}
  \\
\mathit{Tx}_T=(\pi_T,rt,cm_{in},nf,rv,cm_C,\DeltaEpoch)
  \\
\Return(\mathit{voucher},\mathit{note}_C,\mathit{Tx}_T)
\end{gather*}
\end{schemeblock}

\subsubsection{Transfer Relation}

\begin{schemeblock}
\(
  \mathcal{R}_T(\mathbf{x}_T;\mathbf{w}_T)
\)

\begin{gather*}
\Assert pk_S=\Hash(sk_S)
  \,;\,
  \Assert \addrSender=\Hash(pk_S)
  \\
\Assert \addrReceiver=\Hash(pk_R)
  \\
\Assert
  cm_{in}=\CM(\DocumentHash,q_{in},\mathbf{s}_{in},\addrSender,o_{in})
  \\
\Assert\Member(rt,cm_{in},path_{in})=1
  \,;\,
  \Assert nf=\Hash(sk_S,cm_{in})
  \\
\Assert 0<q_T<2^{64}
  \,;\,
  \Assert 0\leq q_C<2^{64}
  \\
\Assert 0\leq s_{T,k},s_{C,k}<2^{64}
  \quad k\in\{1,\ldots,\ell\}
  \\
\Assert q_{in}=q_T+q_C
  \\
\Assert
  q_Cs_{in,k}=q_{in}s_{C,k}+r_{C,k}
  \quad k\in\{1,\ldots,\ell\}
  \\
\Assert 0\leq r_{C,k}<q_{in}
  \quad k\in\{1,\ldots,\ell\}
  \\
\Assert s_{T,k}+s_{C,k}=s_{in,k}
  \quad k\in\{1,\ldots,\ell\}
  \\
\Assert
  rv=
  \Hash(
    \DocumentHash,q_T,\mathbf{s}_T,
    \addrSender,\addrReceiver,\DeltaEpoch,o_{rv}
  )
  \\
\Assert
  cm_C=\CM(\DocumentHash,q_C,\mathbf{s}_C,\addrSender,o_C)
  \\
\end{gather*}
\end{schemeblock}

\subsubsection{\textsc{TransferSC}}

\begin{schemeblock}
\(
  \textsc{TransferSC}(\pi_T,rt,cm_{in},nf,rv,cm_C,\DeltaEpoch)
\)

\begin{gather*}
\mathbf{x}_T=\{rt,cm_{in},nf,rv,cm_C,\DeltaEpoch\}
  \\
\Assert rt\in\List{rt}
  \,;\,
  \Assert nf\notin\List{nf}
  \\
\Assert cm_C\notin\List{cm}
  \,;\,
  \Assert rv\notin\List{rv}
  \\
\Assert
  \Verify(vk_{SC,Transfer},\pi_T,\mathbf{x}_T)=1
  \\
rt_{new}\gets\Append(\MT,cm_C)
  \,;\,
  rvrt_{new}\gets\Append(\rvMT,rv)
  \\
\List{rt}.append(rt_{new})
  \,;\,
  \List{rvrt}.append(rvrt_{new})
  \\
\List{nf}.append(nf)
  \,;\,
  \List{cm}.append(cm_C)
  \,;\,
  \List{rv}.append(rv)
  \\
\Map{rv}[rv]
  \gets
  \CurrentEpoch()+\DeltaEpoch
\end{gather*}
\end{schemeblock}

\subsection{Proceed}

\[
  rv\xrightarrow{\mathrm{Proceed}}cm_{recv}
\]

\noindent\textbf{표기.}
\begin{itemize}[nosep]
  \item $T$: Transfer가 만든 pending portion
  \item $R$: Receiver
  \item $recv$: Receiver가 finalize하여 새로 받는 note
\end{itemize}

\subsubsection{\textsc{ProceedClient}}

\begin{schemeblock}
\(
  \textsc{ProceedClient}(\mathit{voucher},usk_R,upk_R,index_{rv})
\)

\begin{gather*}
  \begin{aligned}
  (&rv,\DocumentHash,q_T,\mathbf{s}_T,
    \addrSender,\addrReceiver,\DeltaEpoch,o_{rv})\\
  &\Leftarrow\mathit{voucher}
  \end{aligned}
  \\
  sk_R\Leftarrow usk_R
  \,;\,
  (\addrReceiver,pk_R)\Leftarrow upk_R
  \\
  (rvrt,path_{rv})
  \gets
  \mathrm{GetPath}_{rvMT}(index_{rv})
  \\
  o_{recv}\sample\mathbb{F}
  \\
  rvnf\gets\Hash(rv,o_{rv})
  \\
  cm_{recv}
  \gets
  \CM(\DocumentHash,q_T,\mathbf{s}_T,\addrReceiver,o_{recv})
  \\
  \mathbf{x}_P=\{rvrt,rv,rvnf,cm_{recv}\}
  \\
  \mathbf{w}_P=\{
    sk_R,pk_R,
    \DocumentHash,q_T,\mathbf{s}_T,
    \addrSender,\addrReceiver,\DeltaEpoch,
    path_{rv},o_{rv},o_{recv}\}
  \\
  \pi_P\gets
  \Prove(ek_{Proceed},\mathbf{x}_P;\mathbf{w}_P)
  \\
  \mathit{note}_{recv}=(cm_{recv},\DocumentHash,q_T,\mathbf{s}_T,o_{recv})
  \\
  \mathit{Tx}_P=(\pi_P,rvrt,rv,rvnf,cm_{recv})
  \\
  \Return(\mathit{note}_{recv},\mathit{Tx}_P)
\end{gather*}
\end{schemeblock}

\subsubsection{Proceed Relation}

\begin{schemeblock}
\(
  \mathcal{R}_P(\mathbf{x}_P;\mathbf{w}_P)
\)

\begin{gather*}
  \Assert pk_R=\Hash(sk_R)
  \,;\,
  \Assert \addrReceiver=\Hash(pk_R)
  \\
  \Assert\Member(rvrt,rv,path_{rv})=1
  \\
  \Assert
  rv=
  \Hash(
    \DocumentHash,q_T,\mathbf{s}_T,
    \addrSender,\addrReceiver,\DeltaEpoch,o_{rv}
  )
  \\
  \Assert rvnf=\Hash(rv,o_{rv})
  \\
  \Assert
  cm_{recv}=\CM(\DocumentHash,q_T,\mathbf{s}_T,\addrReceiver,o_{recv})
\end{gather*}
\end{schemeblock}

\subsubsection{\textsc{ProceedSC}}

\begin{schemeblock}
\(
  \textsc{ProceedSC}(\pi_P,rvrt,rv,rvnf,cm_{recv})
\)

\begin{gather*}
  \mathbf{x}_P=\{rvrt,rv,rvnf,cm_{recv}\}
  \\
  \Assert rvrt\in\List{rvrt}
  \,;\,
  \Assert rv\in\List{rv}
  \\
  \Assert rvnf\notin\List{rvnf}
  \,;\,
  \Assert cm_{recv}\notin\List{cm}
  \\
  \Assert
  \Verify(vk_{SC,Proceed},\pi_P,\mathbf{x}_P)=1
  \\
  rt_{new}\gets\Append(\MT,cm_{recv})
  \\
  \List{rt}.append(rt_{new})
  \,;\,
  \List{cm}.append(cm_{recv})
  \,;\,
  \List{rvnf}.append(rvnf)
\end{gather*}
\end{schemeblock}

\subsection{Recall}

\[
  rv\xrightarrow{\mathrm{Recall}}cm_{return}
\]

\noindent\textbf{표기.}
\begin{itemize}[nosep]
  \item $T$: Transfer가 만든 pending portion
  \item $S$, $R$: 각각 직전 Sender와 Receiver
  \item $return$: Sender에게 되돌아가는 새 finalized note
\end{itemize}

\subsubsection{\textsc{RecallClient}}

\begin{schemeblock}
\(
  \textsc{RecallClient}(\mathit{voucher},usk_S,upk_S,index_{rv})
\)

\begin{gather*}
  \begin{aligned}
  (&rv,\DocumentHash,q_T,\mathbf{s}_T,
    \addrSender,\addrReceiver,\DeltaEpoch,o_{rv})\\
  &\Leftarrow\mathit{voucher}
  \end{aligned}
  \\
  sk_S\Leftarrow usk_S
  \,;\,
  (\addrSender,pk_S)\Leftarrow upk_S
  \\
  (rvrt,path_{rv})
  \gets
  \mathrm{GetPath}_{rvMT}(index_{rv})
  \\
  o_{return}\sample\mathbb{F}
  \\
  rvnf\gets\Hash(rv,o_{rv})
  \\
  cm_{return}
  \gets
  \CM(\DocumentHash,q_T,\mathbf{s}_T,\addrSender,o_{return})
  \\
  \mathbf{x}_R=\{rvrt,rv,rvnf,cm_{return}\}
  \\
  \mathbf{w}_R=\{
    sk_S,pk_S,
    \DocumentHash,q_T,\mathbf{s}_T,
    \addrSender,\addrReceiver,\DeltaEpoch,
    path_{rv},o_{rv},o_{return}\}
  \\
  \pi_R\gets
  \Prove(ek_{Recall},\mathbf{x}_R;\mathbf{w}_R)
  \\
  \mathit{note}_{return}=(cm_{return},\DocumentHash,q_T,\mathbf{s}_T,o_{return})
  \\
  \mathit{Tx}_R=(\pi_R,rvrt,rv,rvnf,cm_{return})
  \\
  \Return(\mathit{note}_{return},\mathit{Tx}_R)
\end{gather*}
\end{schemeblock}

\subsubsection{Recall Relation}

\begin{schemeblock}
\(
  \mathcal{R}_R(\mathbf{x}_R;\mathbf{w}_R)
\)

\begin{gather*}
  \Assert pk_S=\Hash(sk_S)
  \,;\,
  \Assert \addrSender=\Hash(pk_S)
  \\
  \Assert\Member(rvrt,rv,path_{rv})=1
  \\
  \Assert
  rv=
  \Hash(
    \DocumentHash,q_T,\mathbf{s}_T,
    \addrSender,\addrReceiver,\DeltaEpoch,o_{rv}
  )
  \\
  \Assert rvnf=\Hash(rv,o_{rv})
  \\
  \Assert
  cm_{return}=\CM(\DocumentHash,q_T,\mathbf{s}_T,\addrSender,o_{return})
\end{gather*}
\end{schemeblock}

\subsubsection{\textsc{RecallSC}}

\begin{schemeblock}
\(
  \textsc{RecallSC}(\pi_R,rvrt,rv,rvnf,cm_{return})
\)

\begin{gather*}
  \mathbf{x}_R=\{rvrt,rv,rvnf,cm_{return}\}
  \\
  \Assert rvrt\in\List{rvrt}
  \,;\,
  \Assert rv\in\List{rv}
  \\
  \Assert \CurrentEpoch()<\Map{rv}[rv]
  \\
  \Assert rvnf\notin\List{rvnf}
  \,;\,
  \Assert cm_{return}\notin\List{cm}
  \\
  \Assert
  \Verify(vk_{SC,Recall},\pi_R,\mathbf{x}_R)=1
  \\
  rt_{new}\gets\Append(\MT,cm_{return})
  \\
  \List{rt}.append(rt_{new})
  \,;\,
  \List{cm}.append(cm_{return})
  \,;\,
  \List{rvnf}.append(rvnf)
\end{gather*}
\end{schemeblock}

\subsection{Merge}

\[
  \{cm_{old,1},cm_{old,2}\}
  \xrightarrow{\mathrm{Merge}}
  \{cm_{new}\}
\]

\noindent\textbf{표기.}
\begin{itemize}[nosep]
  \item $1,2$: 두 input note의 고정 index
  \item $out$ 또는 $new$: 두 input을 합쳐 만드는 하나의 output note
  \item $k$: State coordinate index
\end{itemize}

\subsubsection{\textsc{MergeClient}}

\begin{schemeblock}
\(
  \textsc{MergeClient}(
    \mathit{note}_1,\mathit{note}_2,usk,upk,
    index_1,index_2)
\)

\begin{gather*}
(cm_i,\DocumentHash_i,q_i,\mathbf{s}_i,o_i)
  \Leftarrow\mathit{note}_i
  \quad i\in\{1,2\}
  \\
sk\Leftarrow usk
  \,;\,
  (addr,pk)\Leftarrow upk
  \\
(rt_i,path_i)\gets\mathrm{GetPath}_{MT}(index_i)
  \quad i\in\{1,2\}
  \\
q_{out}\gets q_1+q_2
  \,;\,
  \mathbf{s}_{out}\gets\mathbf{s}_1+\mathbf{s}_2
  \\
o_{out}\sample\mathbb{F}
  \\
nf_i\gets\Hash(sk,cm_i)
  \quad i\in\{1,2\}
  \\
cm_{out}\gets
  \CM(\DocumentHash_1,q_{out},\mathbf{s}_{out},addr,o_{out})
  \\
\mathbf{x}_M=\{rt_1,rt_2,cm_1,cm_2,nf_1,nf_2,cm_{out}\}
  \\
\mathbf{w}_M=\{
    sk,pk,
    \{\DocumentHash_i,q_i,\mathbf{s}_i\}_{i=1}^{2},
    q_{out},\mathbf{s}_{out},addr,
    \{path_i\}_{i=1}^{2},\{o_i\}_{i=1}^{2},o_{out}\}
  \\
\pi_M\gets
  \Prove(ek_{Merge},\mathbf{x}_M;\mathbf{w}_M)
  \\
\mathit{note}_{out}=(cm_{out},\DocumentHash_1,q_{out},\mathbf{s}_{out},o_{out})
  \\
\mathit{Tx}_M=(\pi_M,rt_1,rt_2,cm_1,cm_2,nf_1,nf_2,cm_{out})
  \\
\Return(\mathit{note}_{out},\mathit{Tx}_M)
\end{gather*}
\end{schemeblock}

\subsubsection{Merge Relation}

\begin{schemeblock}
\(
  \mathcal{R}_M(\mathbf{x}_M;\mathbf{w}_M)
\)

\begin{gather*}
\Assert pk=\Hash(sk)
  \,;\,
  \Assert addr=\Hash(pk)
  \\
\Assert
  cm_i=\CM(\DocumentHash_i,q_i,\mathbf{s}_i,addr,o_i)
  \quad i\in\{1,2\}
  \\
\Assert\Member(rt_i,cm_i,path_i)=1
  \quad i\in\{1,2\}
  \\
\Assert nf_i=\Hash(sk,cm_i)
  \quad i\in\{1,2\}
  \\
\Assert \DocumentHash_1=\DocumentHash_2
  \\
\Assert q_{out}=q_1+q_2
  \,;\,
  \Assert s_{out,k}=s_{1,k}+s_{2,k}
  \quad k\in\{1,\ldots,\ell\}
  \\
\Assert 0\leq q_{out}<2^{64}
  \,;\,
  \Assert 0\leq s_{out,k}<2^{64}
  \quad k\in\{1,\ldots,\ell\}
  \\
\Assert
  cm_{out}=\CM(\DocumentHash_1,q_{out},\mathbf{s}_{out},addr,o_{out})
\end{gather*}
\end{schemeblock}

\subsubsection{\textsc{MergeSC}}

\begin{schemeblock}
\(
  \textsc{MergeSC}(\pi_M,rt_1,rt_2,cm_1,cm_2,nf_1,nf_2,cm_{out})
\)

\begin{gather*}
\mathbf{x}_M=\{rt_1,rt_2,cm_1,cm_2,nf_1,nf_2,cm_{out}\}
  \\
\Assert rt_i\in\List{rt}
  \quad i\in\{1,2\}
  \\
\Assert cm_1\neq cm_2
  \\
\Assert nf_i\notin\List{nf}
  \quad i\in\{1,2\}
  \\
\Assert cm_{out}\notin\List{cm}
  \\
\Assert
  \Verify(vk_{SC,Merge},\pi_M,\mathbf{x}_M)=1
  \\
rt_{new}\gets\Append(\MT,cm_{out})
  \\
\List{nf}.append(nf_1,nf_2)
  \,;\,
  \List{cm}.append(cm_{out})
  \,;\,
  \List{rt}.append(rt_{new})
\end{gather*}
\end{schemeblock}

\subsection{Split}

\[
  \{cm_{old}\}
  \xrightarrow{\mathrm{Split}}
  \{cm_{new,1},cm_{new,2}\}
\]

\noindent\textbf{표기.}
\begin{itemize}[nosep]
  \item $in$: 소비하는 input note
  \item $1,2$: 두 output note의 고정 index
  \item $k$: State coordinate index
  \item Output 2: 내림 계산 대상
  \item Output 1: residual을 받는 대상
\end{itemize}

\subsubsection{\textsc{SplitClient}}

\begin{schemeblock}
\(
  \textsc{SplitClient}(
    \mathit{note}_{in},usk,upk,q_1,q_2,index_{in})
\)

\begin{gather*}
(cm_{in},\DocumentHash,q_{in},\mathbf{s}_{in},o_{in})
  \Leftarrow\mathit{note}_{in}
  \\
sk\Leftarrow usk
  \,;\,
  (addr,pk)\Leftarrow upk
  \\
(rt,path_{in})\gets\mathrm{GetPath}_{MT}(index_{in})
  \\
s_{2,k}\gets
\left\lfloor\frac{q_2s_{in,k}}{q_{in}}\right\rfloor
  \,;\,
r_{2,k}\gets q_2s_{in,k}-q_{in}s_{2,k}
  \quad k\in\{1,\ldots,\ell\}
  \\
\mathbf{s}_1\gets\mathbf{s}_{in}-\mathbf{s}_2
  \\
o_1,o_2\sample\mathbb{F}
  \,;\,
  nf\gets\Hash(sk,cm_{in})
  \\
cm_{out,j}
  \gets
  \CM(\DocumentHash,q_j,\mathbf{s}_j,addr,o_j)
  \quad j\in\{1,2\}
  \\
\mathbf{x}_S=\{rt,cm_{in},nf,cm_{out,1},cm_{out,2}\}
  \\
\mathbf{w}_S=\{
    sk,pk,\DocumentHash,q_{in},\mathbf{s}_{in},
    \{q_j,\mathbf{s}_j\}_{j=1}^{2},addr,\mathbf{r}_2,
    path_{in},o_{in},\{o_j\}_{j=1}^{2}\}
  \\
\pi_S\gets
  \Prove(ek_{Split},\mathbf{x}_S;\mathbf{w}_S)
  \\
\mathit{note}_{out,j}=(cm_{out,j},\DocumentHash,q_j,\mathbf{s}_j,o_j)
  \quad j\in\{1,2\}
  \\
\mathit{Tx}_S=(\pi_S,rt,cm_{in},nf,cm_{out,1},cm_{out,2})
  \\
\Return(\mathit{note}_{out,1},\mathit{note}_{out,2},\mathit{Tx}_S)
\end{gather*}
\end{schemeblock}

\subsubsection{Split Relation}

\begin{schemeblock}
\(
  \mathcal{R}_S(\mathbf{x}_S;\mathbf{w}_S)
\)

\begin{gather*}
\Assert pk=\Hash(sk)
  \,;\,
  \Assert addr=\Hash(pk)
  \\
\Assert
  cm_{in}=\CM(\DocumentHash,q_{in},\mathbf{s}_{in},addr,o_{in})
  \\
\Assert\Member(rt,cm_{in},path_{in})=1
  \,;\,
  \Assert nf=\Hash(sk,cm_{in})
  \\
\Assert 0\leq q_1,q_2<2^{64}
  \\
\Assert 0\leq s_{1,k},s_{2,k}<2^{64}
  \quad k\in\{1,\ldots,\ell\}
  \\
\Assert q_{in}=q_1+q_2
  \\
\Assert
  q_2s_{in,k}=q_{in}s_{2,k}+r_{2,k}
  \quad k\in\{1,\ldots,\ell\}
  \\
\Assert 0\leq r_{2,k}<q_{in}
  \quad k\in\{1,\ldots,\ell\}
  \\
\Assert s_{1,k}+s_{2,k}=s_{in,k}
  \quad k\in\{1,\ldots,\ell\}
  \\
\Assert
  cm_{out,j}=\CM(\DocumentHash,q_j,\mathbf{s}_j,addr,o_j)
  \quad j\in\{1,2\}
\end{gather*}
\end{schemeblock}

Output quantity가 0이면 해당 output의 모든 State coordinate도 0으로 강제된다.
$q_2=0$이면 floor/remainder relation에서 $s_{2,k}=r_{2,k}=0$이고,
$q_1=0$이면 $q_2=q_{in}$이므로 $s_{2,k}=s_{in,k}$와 $s_{1,k}=0$이 성립한다.
$q_{in}=0$인 note는 $0\leq r_{2,k}<q_{in}$을 만족할 수 없어 Split할 수 없다.

\subsubsection{\textsc{SplitSC}}

\begin{schemeblock}
\(
  \textsc{SplitSC}(\pi_S,rt,cm_{in},nf,cm_{out,1},cm_{out,2})
\)

\begin{gather*}
\mathbf{x}_S=\{rt,cm_{in},nf,cm_{out,1},cm_{out,2}\}
  \\
\Assert rt\in\List{rt}
  \,;\,
  \Assert nf\notin\List{nf}
  \\
\Assert
  \Verify(vk_{SC,Split},\pi_S,\mathbf{x}_S)=1
  \\
\List{nf}.append(nf)
  \\
\begin{aligned}
&\Assert cm_{out,j}\notin\List{cm}
  \,;\,
  \List{cm}.append(cm_{out,j})\\
&rt_j\gets\Append(\MT,cm_{out,j})
  \,;\,
  \List{rt}.append(rt_j)
\end{aligned}
  \quad j=1,2\ \text{ 순서로 실행}
\end{gather*}
\end{schemeblock}

\subsection{Process}

현재 Process는 arbitrary industrial process를 일반화한 Relation이 아니라, 호출자가
공개한 $P$와 $A$에 대한 generic State Allocation Relation이다. 이번 POC는 Recipe,
MaterialClassId와 Quantity transition을 검증하지 않는다. 따라서 Process proof가
보장하는 범위는 input State가 공개된 $P,A$에 따라 output State로 일관되게 전이되었다는
사실이며, $P,A$ 자체가 실제 공정에 맞는지는 보장하지 않는다.

\[
  1\leq m\leq M_{\max}=a,
  \qquad
  1\leq n\leq N_{\max}=a,
  \qquad
  a\in\{1,2,3\}
\]

기능 확인용 기본 POC는 $a=3$이다. Process arity benchmark는 $d=32,\ell=3$으로
고정하고 $a=1,2,3$인 circuit을 각각 별도로 compile한다. 기준 scenario에서 모든 slot을 활성화해
$m=n=a$인 1-to-1, 2-to-2와 3-to-3만 측정한다. 비대칭 arity는 benchmark에서 제외한다.

첫 구현 milestone의 Process circuit은 selector를 아직 포함하지 않고 $m=n=3$을 직접
강제한다. 아래 $m,n$ 기반 active-slot 정의는 3-to-3 vertical slice가 완료된 뒤 추가하는
일반 최대 3-slot profile의 요구사항이다.

\[
  D:=10^9
\]

Public $m,n$은 각각 active input/output의 개수를 나타낸다. Circuit은 index
$1,\ldots,m$과 $1,\ldots,n$에 해당하는 fixed slot만 active relation에 포함한다.
별도 $e^{in},e^{out}$ vector는 statement와 witness에 포함하지 않는다. Inactive input의
$rt_i,cm_i,nf_i$와 inactive output의 $cm_j$는 0이며, inactive output의 $A$ row도 0이다.
고정 circuit에서 $m,n$을 slot별 constraint에 적용하는 algebraic gating은 Development
Specification에서 정의한다.

\[
  \{cm_{old,i}\}_{i=1}^{m}
  \xrightarrow{\mathrm{Process}}
  \{cm_{new,j}\}_{j=1}^{n}
\]

\noindent\textbf{표기.}
\begin{itemize}[nosep]
  \item $i$, $j$: 각각 input index와 output index
  \item $k$, $\ell$: 각각 State coordinate index와 StateVector 길이
  \item $[t]:=\{1,\ldots,t\}$: index 범위를 짧게 쓰기 위한 표기
  \item $m$, $n$: 실제 input 개수와 output 개수
  \item $a$: circuit의 maximum arity
  \item $\mathbf{s}_{in,i}$: $i$번째 input note의 StateVector
  \item $\mathbf{s}_{out,j}$: $j$번째 output note의 StateVector
  \item $\mathbf{s}_{in}:=\sum_{i=1}^{m}\mathbf{s}_{in,i}$: input aggregate StateVector
  \item $\mathbf{s}_{out}:=\sum_{j=1}^{n}\mathbf{s}_{out,j}$: output aggregate StateVector
  \item $P$: 공정 변화량
  \item $A\in\mathbb{F}^{n\times\ell}$: Client와 transaction이 전달하는 active
        StateAllocationMatrix
  \item $A^{pad}\in\mathbb{F}^{N_{\max}\times\ell}$: $A$ 뒤에 zero row를 붙여
        verifier statement에 사용하는 fixed-size matrix
\end{itemize}

\Needspace{0.68\textheight}
\subsubsection{\textsc{ProcessClient}}

$P$는 Client가 실제 공정 결과를 나타내기 위해 공개하는 semantic process delta다.
Benchmark profile의 State 길이에 따라 다음 prefix를 사용한다.

\[
\begin{array}{c|c|c}
\ell & P & P_{pub} \\
\hline
0 & () & () \\
1 & (-p_{totalLoss}) & (p_{totalLoss}) \\
2 & (-p_{totalLoss},+p_{carbonAdded}) & (p_{totalLoss},p_{carbonAdded}) \\
3 & (-p_{totalLoss},+p_{carbonAdded},+p_{recycledDelta}) &
(p_{totalLoss},p_{carbonAdded},p_{recycledDelta})
\end{array}
\]

Transaction과 public input에는 음수 Field Element 대신 $P_{pub}$의 unsigned 값을
포함한다. $\ell=3$에서는 recycled coordinate도 생략하지 않는다. Relation은
$\mathbf{s}_{out}=\mathbf{s}_{in}+P$와 recycled aggregate 보존을 함께 검증하므로
유효한 proof에서 $p_{recycledDelta}=0$이 도출된다.
$\mathbf{s}_{in}$과 $\mathbf{s}_{out}$은 전달하지 않으며 Relation이 각각 input과
output StateVector의 합으로 계산한다.
아래 수식은 active slot만 표시한다. Client는 $\Prove$ 호출 전에 Relation의 inactive
slot 규칙에 따라 fixed-size statement와 witness의 나머지 값을 0으로 채운다.

\begin{schemeblock}
\(
  \textsc{ProcessClient}(
    \{\mathit{note}_{in,i},index_i\}_{i=1}^{m},usk,upk,
    \{\DocumentHash_{out,j},q_{out,j}\}_{j=1}^{n},
    P_{pub},A)
\)

\begin{gather*}
\text{\bfseries 1. Input note, 소유자 key와 MT path 준비}
  \\
(cm_{in,i},\DocumentHash_{in,i},q_{in,i},\mathbf{s}_{in,i},o_{in,i})
  \Leftarrow\mathit{note}_{in,i}
  \quad i\in[m]
  \\
sk\Leftarrow usk
  \,;\,
  (addr,pk)\Leftarrow upk
  \\
(rt_i,path_i)\gets\mathrm{GetPath}_{MT}(index_i)
  \quad i\in[m]
  \\
\text{\bfseries 2. Process delta 반영과 output State 계산}
  \\
P\gets(-p_{totalLoss},+p_{carbonAdded},+p_{recycledDelta})
  \quad (\ell=3)
  \\
\mathbf{s}_{in}\gets\sum_{i=1}^{m}\mathbf{s}_{in,i}
  \,;\,
\mathbf{s}_{out}\gets\mathbf{s}_{in}+P
  \\
s_{out,j,k}\gets
\left\lfloor\frac{s_{out,k}A_{j,k}}{D}\right\rfloor
  \,;\,
r_{j,k}\gets s_{out,k}A_{j,k}-Ds_{out,j,k}
  \quad j\in[n]\setminus\{1\},\;k\in[\ell]
  \\
s_{out,1,k}\gets s_{out,k}-\sum_{j=2}^{n}s_{out,j,k}
  \quad k\in[\ell]
  \\
\text{\bfseries 3. Nullifier와 active output note 생성}
  \\
nf_i\gets\Hash(sk,cm_{in,i})
  \quad i\in[m]
  \\
o_{out,j}\sample\mathbb{F}
  \quad j\in[n]
  \\
cm_{out,j}
  \gets
  \CM(\DocumentHash_{out,j},q_{out,j},\mathbf{s}_{out,j},addr,o_{out,j})
  \quad j\in[n]
  \\
\mathit{note}_{out,j}=(
  cm_{out,j},\DocumentHash_{out,j},q_{out,j},
  \mathbf{s}_{out,j},o_{out,j})
  \quad j\in[n]
  \\
\text{\bfseries 4. Fixed-size statement와 witness 구성}
  \\
A^{pad}\gets\operatorname{PadZeroRows}_{N_{\max}}(A)
  \\
\mathbf{x}_{Proc}=\{
    m,n,P_{pub},A^{pad},
    \{rt_i,cm_{in,i},nf_i\}_{i=1}^{M_{\max}},
    \{cm_{out,j}\}_{j=1}^{N_{\max}}\}
  \\
\begin{aligned}
  \mathbf{w}_{Proc}=\{
    &sk,pk,\\
    &\{\DocumentHash_{in,i},q_{in,i},\mathbf{s}_{in,i}\}_{i=1}^{M_{\max}},\\
    &\{\DocumentHash_{out,j},q_{out,j},\mathbf{s}_{out,j}\}_{j=1}^{N_{\max}},\\
    &addr,\{r_{j,k}\}_{j=2,k=1}^{N_{\max},\ell},
    \{path_i\}_{i=1}^{M_{\max}},\\
    &\{o_{in,i}\}_{i=1}^{M_{\max}},\{o_{out,j}\}_{j=1}^{N_{\max}}\}
  \end{aligned}
  \\
\text{\bfseries 5. Proof와 active-size transaction 생성}
  \\
\pi_{Proc}
  \gets
  \Prove(ek_{Process},\mathbf{x}_{Proc};\mathbf{w}_{Proc})
  \\
\mathit{Tx}_{Proc}=(
    \pi_{Proc},P_{pub},A,
    \{rt_i,cm_{in,i},nf_i\}_{i=1}^{m},
    \{cm_{out,j}\}_{j=1}^{n})
  \\
\Return(\{\mathit{note}_{out,j}\}_{j=1}^{n},\mathit{Tx}_{Proc})
\end{gather*}
\end{schemeblock}

\Needspace{0.16\textheight}
\subsubsection{Process Relation}

아래 본문은 기능 확인용 $\ell=3$ profile을 기준으로 한다. $\ell<3$인 benchmark
profile은 StateVector prefix만 사용하고 존재하지 않는 coordinate의 delta relation을
compile하지 않는다. 따라서 $\ell$은 아래 Relation의 runtime 분기값이 아니다.

\begin{schemeblock}
\(
  \mathcal{R}_{Proc}
  (\mathbf{x}_{Proc};\mathbf{w}_{Proc})
\)

\begin{gather*}
\text{\bfseries 1. 소유자와 active arity}
  \\
\Assert pk=\Hash(sk)
  \,;\,
  \Assert addr=\Hash(pk)
  \\
\Assert 1\leq m\leq M_{\max}
  \,;\,
\Assert 1\leq n\leq N_{\max}
  \\
\text{\bfseries 2. Active input의 opening, membership와 nullifier}
  \\
\Assert
  cm_{in,i}=\CM(
    \DocumentHash_{in,i},q_{in,i},\mathbf{s}_{in,i},addr,o_{in,i})
  \quad i\in[m]
  \\
\Assert\Member(rt_i,cm_{in,i},path_i)=1
  \,;\,
  \Assert nf_i=\Hash(sk,cm_{in,i})
  \quad i\in[m]
  \\
\text{\bfseries 3. State aggregate와 active output range}
  \\
\mathbf{s}_{in}:=\sum_{i=1}^{m}\mathbf{s}_{in,i}
  \\
\mathbf{s}_{out}:=\sum_{j=1}^{n}\mathbf{s}_{out,j}
  \\
\Assert 1\leq q_{out,j}<2^{64}
  \,;\,
  \Assert 0\leq s_{out,j,k}<2^{64}
  \quad j\in[n],\;k\in[\ell]
  \\
\text{\bfseries 4. Public process delta}
  \\
P:=(-p_{totalLoss},+p_{carbonAdded},+p_{recycledDelta})
  \\
\Assert 0\leq p_{totalLoss}<2^{66}
  \quad\bigl(M_{\max}=3:\;3(2^{64}-1)<2^{66}\bigr)
  \\
\Assert 0\leq p_{carbonAdded}<2^{66}
  \quad\bigl(N_{\max}=3:\;3(2^{64}-1)<2^{66}\bigr)
  \\
  \Assert \mathbf{s}_{out}=\mathbf{s}_{in}+P
  \,;\,
  \Assert s_{out,recycled}=s_{in,recycled}
  \\
\text{\bfseries 5. State allocation과 output 1 residual}
  \\
\Assert 0\leq A^{pad}_{j,k}\leq D
  \quad j\in[n],\;k\in[\ell]
  \\
\Assert s_{out,k}A^{pad}_{j,k}=Ds_{out,j,k}+r_{j,k}
  \quad j\in[n]\setminus\{1\},\;k\in[\ell]
  \\
\Assert 0\leq r_{j,k}<D
  \quad j\in[n]\setminus\{1\},\;k\in[\ell]
  \\
\Assert\sum_{j=1}^{n}A^{pad}_{j,k}=D
  \quad k\in[\ell]
  \\
\text{\bfseries 6. Inactive slot의 zero padding}
  \\
\Assert rt_i=cm_{in,i}=nf_i=0
  \quad i\in[M_{\max}]\setminus[m]
  \\
\Assert
  \DocumentHash_{in,i}=q_{in,i}=s_{in,i,k}=o_{in,i}=0
  \quad i\in[M_{\max}]\setminus[m],\;k\in[\ell]
  \\
\Assert path_i=\mathbf{0}^{d}
  \quad i\in[M_{\max}]\setminus[m]
  \\
\Assert cm_{out,j}=0
  \,;\,
  \Assert A^{pad}_{j,k}=0
  \quad j\in[N_{\max}]\setminus[n],\;k\in[\ell]
  \\
\Assert
  \DocumentHash_{out,j}=q_{out,j}=s_{out,j,k}=o_{out,j}=r_{j,k}=0
  \quad j\in[N_{\max}]\setminus[n],\;k\in[\ell]
  \\
\text{\bfseries 7. Active output commitment binding}
  \\
\Assert
  cm_{out,j}=\CM(
    \DocumentHash_{out,j},q_{out,j},\mathbf{s}_{out,j},addr,o_{out,j})
  \quad j\in[n]
\end{gather*}
\end{schemeblock}

\Needspace{0.38\textheight}
\subsubsection{\textsc{ProcessSC}}

\begin{schemeblock}
\(
  \textsc{ProcessSC}(\mathit{Tx}_{Proc})
\)

\begin{gather*}
\text{\bfseries 1. Active array 길이와 public input shape 확인}
  \\
 m\gets\operatorname{len}(\{cm_{in,i}\})
  \,;\,
  n\gets\operatorname{len}(\{cm_{out,j}\})
  \\
\Assert m=\operatorname{len}(\{rt_i\})=\operatorname{len}(\{nf_i\})
  \,;\,
  \Assert 1\leq m\leq M_{\max}
  \,;\,
  \Assert 1\leq n\leq N_{\max}
  \\
\Assert \operatorname{len}(P_{pub})=\ell
  \,;\,
  \Assert \operatorname{len}(\operatorname{vec}(A))=n\ell
  \,;\,
  A^{pad}\gets\operatorname{PadZeroRows}_{N_{\max}}(A)
  \\
\text{\bfseries 2. Fixed-size verifier public input 구성}
  \\
\{rt_i,cm_{in,i},nf_i\}_{i=1}^{M_{\max}}
  \gets\operatorname{PadZero}_{M_{\max}}(\{rt_i,cm_{in,i},nf_i\}_{i=1}^{m})
  \\
\{cm_{out,j}\}_{j=1}^{N_{\max}}
  \gets\operatorname{PadZero}_{N_{\max}}(\{cm_{out,j}\}_{j=1}^{n})
  \\
\mathbf{x}_{Proc}=\{
    m,n,P_{pub},A^{pad},
    \{rt_i,cm_{in,i},nf_i\}_{i=1}^{M_{\max}},
    \{cm_{out,j}\}_{j=1}^{N_{\max}}\}
  \\
\text{\bfseries 3. Proof 이전 root, input 중복과 spent 검사}
  \\
\Assert rt_i\in\List{rt}
  \quad i\in[m]
  \\
\Assert cm_{in,i}\neq cm_{in,i'}
  \quad 1\leq i<i'\leq m
  \\
\Assert nf_i\notin\List{nf}
  \quad i\in[m]
  \\
\text{\bfseries 4. Process proof 검증}
  \\
\Assert
  \Verify(
    vk_{SC,Process},
    \pi_{Proc},\mathbf{x}_{Proc})=1
  \\
\text{\bfseries 5. Nullifier와 output commitment 상태 갱신}
  \\
\List{nf}.append(\{nf_i\}_{i=1}^{m})
  \\
\begin{aligned}
&\Assert cm_{out,j}\notin\List{cm}
  \,;\,
  \List{cm}.append(cm_{out,j})\\
&rt_{new,j}\gets\Append(\MT,cm_{out,j})
  \,;\,
  \List{rt}.append(rt_{new,j})
\end{aligned}
  \quad j=1,\ldots,n\ \text{ 순서로 실행}
\end{gather*}
\end{schemeblock}

\subsection{Exit}

\[
  \{cm_{old}\}
  \xrightarrow{\mathrm{Exit}}
  \varnothing
\]

\noindent\textbf{표기.}
\begin{itemize}[nosep]
  \item $in$ 또는 $old$: 공급망에서 제거하기 위해 소비하는 finalized note
\end{itemize}

\Needspace{0.38\textheight}
\subsubsection{\textsc{ExitClient}}

\begin{schemeblock}
\(
  \textsc{ExitClient}(\mathit{note}_{in},index_{in},usk,upk)
\)

\begin{gather*}
\text{\bfseries 1. Input note, 소유자 key와 MT path 준비}
  \\
(cm_{in},\DocumentHash,q_{in},\mathbf{s}_{in},o_{in})
  \Leftarrow\mathit{note}_{in}
  \\
sk\Leftarrow usk
  \,;\,
  (addr,pk)\Leftarrow upk
  \\
(rt,path_{in})\gets\mathrm{GetPath}_{MT}(index_{in})
  \\
\text{\bfseries 2. Nullifier와 proof input 구성}
  \\
nf\gets\Hash(sk,cm_{in})
  \\
\mathbf{x}_{Exit}=\{rt,cm_{in},nf\}
  \\
\mathbf{w}_{Exit}=\{
    sk,pk,\DocumentHash,q_{in},\mathbf{s}_{in},addr,path_{in},o_{in}\}
  \\
\text{\bfseries 3. Proof와 transaction 생성}
  \\
\pi_{Exit}\gets
  \Prove(ek_{Exit},\mathbf{x}_{Exit};\mathbf{w}_{Exit})
  \\
\mathit{Tx}_{Exit}=(\pi_{Exit},rt,cm_{in},nf)
  \\
\Return\mathit{Tx}_{Exit}
\end{gather*}
\end{schemeblock}

\subsubsection{Exit Relation}

\begin{schemeblock}
\(
  \mathcal{R}_{Exit}(\mathbf{x}_{Exit};\mathbf{w}_{Exit})
\)

\begin{gather*}
\text{\bfseries 1. 소유자와 input commitment opening}
  \\
\Assert pk=\Hash(sk)
  \,;\,
  \Assert addr=\Hash(pk)
  \\
\Assert
  cm_{in}=\CM(\DocumentHash,q_{in},\mathbf{s}_{in},addr,o_{in})
  \\
\text{\bfseries 2. MT membership와 nullifier binding}
  \\
\Assert\Member(rt,cm_{in},path_{in})=1
  \,;\,
  \Assert nf=\Hash(sk,cm_{in})
\end{gather*}
\end{schemeblock}

\subsubsection{\textsc{ExitSC}}

\begin{schemeblock}
\(
  \textsc{ExitSC}(\mathit{Tx}_{Exit})
\)

\begin{gather*}
\text{\bfseries 1. Public input과 proof 이전 상태 검사}
  \\
\mathbf{x}_{Exit}=\{rt,cm_{in},nf\}
  \\
\Assert rt\in\List{rt}
  \,;\,
  \Assert nf\notin\List{nf}
  \\
\text{\bfseries 2. Exit proof 검증}
  \\
\Assert
  \Verify(vk_{SC,Exit},\pi_{Exit},\mathbf{x}_{Exit})=1
  \\
\text{\bfseries 3. Nullifier spent 처리}
  \\
\List{nf}.append(nf)
\end{gather*}
\end{schemeblock}

\end{document}
